\clearpage
\fancyhead[L]{\small MOON STARS.}
\hypertarget{native-448-mondsterne}{}
\begin{center}
{\large MOON STARS.}
\end{center}

\begin{center}
Astronomische Nachrichten. Vol. II. No. 26. Pages 17..20. February 1823.
\end{center}

\begin{center}
Calculating the difference in longitude between two places from corresponding culminations of the Moon.\\
By Professor \textsc{Nicolai}, Director of the Observatory in Mannheim [to \textsc{Schumacher.}].
\end{center}

In these lines I shall try today to answer your kind request that I tell you how I calculate the difference in longitude from corresponding culminations of the Moon.

As is well known, this method rests on observing at two places the true difference in right ascension between one limb of the Moon and some fixed stars agreed on beforehand. Let the difference between the true right ascensions of the Moon's limb and a fixed star (taken as positive when the Moon's right ascension exceeds the star's) be $=t$ at the western observing site and $=\tau$ at the eastern one. Then $t-\tau$ will be the true difference in right ascension of the Moon's limb over the time elapsed between the two observations. If several fixed stars have been compared at both places, one takes the arithmetic mean of the various values of $t-\tau$, provided the observations at the two sites are equally precise. We shall see below how to proceed when this condition fails with respect to the number of wires observed. --- What we evidently need, however, is the true difference in right ascension of the \emph{centre} of the Moon. It will differ somewhat from the one just obtained, because the Moon's semidiameter in R.A. changes between observations, partly through a change in the true semidiameter itself and partly through a change in the Moon's declination. For the time of the Moon's culmination at the western observing site, let $r$ be its \emph{true} semidiameter (as seen from the centre of the Earth) and $d$ its true declination, and for the time of culmination at the eastern site, denote the same quantities by $\rho$ and $\delta$. Then

\[
t-\tau \mp \tfrac{1}{15}\left(\frac{r}{\cos d}-\frac{\rho}{\cos\delta}\right) = \Delta
\]
\noindent gives the true difference in right ascension of the centre of the Moon for the two instants of observation. The lower sign in this expression applies when the western or first limb of the Moon is observed, the upper sign when the eastern or second limb is observed. In all cases where the difference in longitude between the two observing sites is not very large, this correction can be neglected altogether. It is also
\SourceLeaf{449}%
easy to see that the absolute values of $r$ and $d$ need not be known to great precision, since only the \emph{change} in these quantities over the interval between observations matters here.

To obtain the difference in longitude between the two observing sites from $\Delta$, we need only one element, namely the hourly motion of the Moon in R.A. If we denote this by $h$ and the difference in longitude in time by $x$, then

\[
x = \left(\frac{1}{h}\,m\cdot 15\cdot 3600 - 1\right)\Delta = n\Delta
\]
\noindent when we set the factor $\frac{1}{h}m\cdot15\cdot3600-1=n$. In this expression the factor $m$ depends on the kind of time for which $h$ is calculated. If $h$ is the motion in one sidereal hour, then $m=1$. If $h$ is the motion in one mean solar hour, $m$ denotes the ratio of the mean solar day to the sidereal day ($=1.00274$). Finally, if $h$ is the motion in one apparent solar hour, $m$ denotes the corresponding ratio of the apparent solar day to the sidereal day. --- As for calculating the hourly motion $h$, it is worth noting that there would be no advantage in deriving it from the lunar tables by directly calculating two positions of the Moon just one hour apart. For although the lunar tables are calculated to tenths of a second, the large number of equations to be summed means that one can never vouch for a lunar position calculated from them to within one second. It is therefore clear that the hourly motion calculated in this way would often be appreciably in error. It is much more reliable, and indeed highly precise, to calculate several positions of the Moon at \emph{considerable} intervals, for example every 12 hours, and to derive the hourly motion from them by \emph{precise interpolation}. This will distribute evenly the small errors that still affect the individual positions. For this operation one can safely use the positions of the Moon calculated every 12 hours in the ephemerides. --- Given the \emph{variability} of the hourly motion itself, one must also take care to calculate the $h$ used to determine the difference in longitude \emph{for the mean of the two observation times reduced to the same meridian}. Let $A$ be the right ascension of the Moon for the hour $\theta$ at a given meridian, and $A'$ the same for the hour $\theta+1$ at the same meridian. Then $A'-A$ or $h$ will be the hourly motion of the Moon in right ascension for the time $\theta+30^{\rm m}$. One therefore calculates several successive hourly motions for the meridian of the tables or ephemerides, reduces both observation times to this meridian of the tables or ephemerides by means of the approximately known difference in longitude, and obtains the $h$ to be used for the mean of those times by interpolation from the hourly motions.

Since the comparison stars and the Moon cannot always be observed at every wire of the meridian telescope on every evening and by every observer, it is clear that the results cannot all be assigned the same precision. This too must be taken into account. Court Councillor \textsc{Gauss} was kind enough to send me the expressions for the relative precision of the individual results some time ago. They follow readily from the method of least squares. Suppose the Moon was observed at $l$ wires at one place and at $l'$ at the other, and the various stars at $a, b, c$ etc. wires at one place and at $a', b', c'$ etc. at the other. Set

\[
\frac{l\,l'}{l+l'}=\lambda,\qquad \frac{a\,a'}{a+a'}=\alpha,\qquad \frac{b\,b'}{b+b'}=\ObservedBeta,\qquad \frac{c\,c'}{c+c'}=\gamma,
\]
\noindent etc. Now the mean of the results from the individual comparison stars must be taken by assigning them the respective weights (the square of the precision) $\alpha, \ObservedBeta, \gamma$ etc. The result thus obtained
\SourceLeaf{450}%
for $t-\tau$ or for $\Delta$ then has the weight $\dfrac{\sigma\lambda}{\sigma+\lambda}$ if we set $\alpha+\ObservedBeta+\gamma+\text{etc.}=\sigma$. If the coefficient by which $\Delta$ is multiplied to obtain the difference in longitude is denoted by $n$, as above, the weight of the result for that difference is $=\dfrac{\sigma\lambda}{(\sigma+\lambda)nn}$ and the weight of the result from all the determinations is $=\Sigma\cdot\dfrac{\sigma\lambda}{(\sigma+\lambda)nn}$. Or, if the probable error of a wire crossing is set $=\varepsilon$, the probable error of the final result is

\[
=\frac{\varepsilon}{\BareRoot{\Sigma\,\dfrac{\sigma\lambda}{(\sigma+\lambda)nn}}}
\]

As is evident, this calculation of the probable error assumes that all observed crossings are assigned equal reliability, without distinguishing between the Moon and the stars, between the various declinations, or between observers. It would be easy to account for all these circumstances if the corresponding coefficients were known. All this, however, is probably of little consequence. As for the different declinations, one can safely regard all crossings observed within the Moon's zodiac as equally precise, especially since the circumstances of each observation, the state of the air, the lower or higher position of the Moon and stars, and so on, will have far more effect on the precision of the crossings here than greater or lesser declinations. It will, however, be advisable to use for $s$ in the final formula above a \emph{mean} value from all the crossings observed within the Moon's zodiac.

\begin{center}
{\small ---\,---\,---\,---\,---\,---\,---\,---\,---\,---\,---\,---\,---\,---\,---\,---\,---\,---\,---\,---\,---\,---}
\end{center}

As you can see, this whole calculation is extremely simple and, since the Conn.\ de Tems has spared us the trouble of calculating the Moon's right ascensions ourselves, also very convenient and brief. --- In closing, I would add that, \emph{in full rigour}, the value of $x+\Delta$ is actually obtained from $\Delta$ by a series in successive odd powers of $\Delta$, whose first term is $\dfrac{1}{h}m\cdot15\cdot3600\Delta$ (with $h$ taken for the mean of the two observation times). The second term, which depends on $\Delta^3$, is already so small, however, that it amounts to about $0^{\mathrm{s}}1$ only for a difference in longitude of two to two and a half hours.
For a \emph{very large} difference in longitude, however, I would prefer an \emph{indirect} procedure. I would calculate the difference in the Moon's right ascensions corresponding to an \emph{assumed} difference in longitude, compare this calculated difference with the one actually observed, and use their discrepancy and the hourly motion to determine the correction still to be applied to the assumed difference in longitude.

\begin{flushright}
\textsc{Nicolai.}
\end{flushright}
